\loai{Tính công suất hao phí}
\begin{vidu} %[3P3-15.1B][loai1] 
	Một nhà máy điện sinh ra một công suất 100000 kW và cần truyền tải tới nơi tiêu thụ. Biết hiệu suất truyền tải là 90\%. Công suất hao phí trên đường truyền là
	\choice
	{\True 10000 kW}
	{1000 kW}
	{100 kW}
	{10 kW}
	\loigiai{
		Hiệu suất truyền tải: $ H=1-\dfrac{P_{hp}}{P}=90\%\Rightarrow \dfrac{P_{hp}}{P}=0,1\Rightarrow {{P}_{hp}}=0,1. 100000=10000\ kW$.
	}
\end{vidu}
\loai{Tính điện trở}
\begin{vidu} %[3P3-15.1B][loai1] 
	Truyền một công suất 500 kW từ một trạm phát điện đến nơi tiêu thụ bằng đường dây một pha. Biết công suất hao phí trên đường dây là 10 kW, điện áp hiệu dụng ở trạm phát là 35 kV. Coi hệ số công suất của mạch truyền tải điện bằng 1. Điện trở tổng cộng của đường dây tải điện là
	\choice
	{$55\ \Omega$}
	{\True $49\ \Omega$}
	{$38\ \Omega$}
	{$52\ \Omega$}
	\loigiai{
		Công suất hao phí trên đường dây tải điện là: $\Delta P=\left( \dfrac{P}{U\cos \varphi} \right)^2. R\Rightarrow R=\dfrac{\Delta P. {U^2}}{{{P}^2}}=\dfrac{{{10000. 35000}^2}}{{{500000}^2}}=49\ \Omega$.
	}
\end{vidu}
\begin{vidu} %[3P3-15.4K][loai1] 
	Người ta cần truyền tải điện năng từ máy hạ thế có điện áp đầu ra 200 V đến một hộ gia đình cách 1 km. Công suất tiêu thụ ở đầu ra của máy biến áp cho hộ gia đình đó là 10 kW và yêu cầu độ giảm điện áp trên dây không quá 20 V. Điện trở suất dây dẫn là $\rho =2,8.10^{-8}$ và tải tiêu thụ là điện trở. Tiết diện dây dẫn phải thoả mãn
	\choice
	{\True $S = 1,4\ cm^2$}
	{$S = 0,7\ cm^2$}
	{$S = 0,7\ m^2$}
	{$S = 1,4\ m^2$}
	\loigiai{
		\begin{itemize}
			\item Ta có: $I = \dfrac{P}{U} = \dfrac{10000}{200} = 50\ A.$
			\item Độ giảm điện thế không quá 20 V thì $R\le \dfrac{20}{50}=0,4\ \Omega $.
			\item Lại có điện trở $R=\dfrac{\rho \ell}{S}\Rightarrow R\le 0,4\ \Omega $ thì $S\ge \dfrac{\rho \ell}{R}=\dfrac{2,8.10^{-8}.2.1000}{0,4}=1,4.10^{-4}\ m^2=1,4\ cm^2$.
		\end{itemize}
	}
\end{vidu}
\loai{Tính công suất tiêu thụ}
\begin{vidu} %[3P3-15.1K][loai1] 
	Điện năng được truyền từ một máy biến áp ở A, ở nhà máy điện tới một máy hạ áp ở nơi tiêu thụ bằng hai dây đồng có điện trở tổng cộng là $40\ \Omega$. Cường độ dòng điện trên đường dây tải là 50 A. Công suất tiêu hao trên đường dây tải bằng 5\% công suất tiêu thụ ở B. Công suất tiêu thụ ở B bằng?
	\choice
	{200 kW}
	{\True 2 MW}
	{2 kW}
	{200 W}
	\loigiai{
		\begin{itemize}
			\item Gọi công suất tiêu thụ ở B là P.
			\item Ta có: $\Delta P=I^2R=0,05P\Rightarrow $ $P=\dfrac{I^2R}{0,05}=\dfrac{{{50}^2}. 40}{0,05}={{2. 10}^{6}}\ W=2\ MW$.
		\end{itemize}
	}
\end{vidu}
\loai{Tính tỉ số số vòng dây}
\begin{vidu} %[3P3-15.1B][loai1] 
	Khi truyền tải điện năng có công suất không đổi đi xa với đường dây tải điện một pha có điện trở R xác định. Để công suất hao phí trên đường dây tải điện giảm đi 100 lần thì ở nơi truyền đi phải dùng một máy biến áp lí tưởng có tỉ số vòng dây giữa cuộn thứ cấp và cuộn sơ cấp là
	\choice
	{100}
	{\True 10}
	{50}
	{40}
	\loigiai{
		\begin{itemize}
			\item Công suất hao phí: $\Delta P=\dfrac{{{P}^2} R}{{U^2}}$.
			\item Để công suất hao phí giảm 100 lần thì điện áp phải giảm 10 lần. 
			\item Do đó tỉ số vòng dây giữa cuộn thứ cấp và cuộn sơ cấp là 10.
	\end{itemize}}
\end{vidu}
\loai{Tính điện năng hao phí}
\begin{vidu} %[3P3-15.1B][loai1] 
	Ở trạm phát điện xoay chiều một pha có điện áp hiệu dụng 110 kV, truyền đi công suất điện 1000 kW trên đường dây dẫn có điện trở $20\ \Omega$. Hệ số công suất của đoạn mạch $\cos \varphi =0,9$. Điện năng hao phí trên đường dây trong 30 ngày là
	\choice
	{5289 kWh}
	{61,2 kWh}
	{145,5 kWh}
	{\True 1469 kWh}
	\loigiai{
		\begin{itemize}
			\item Công suất hao phí trên đường dây tải điện là: $\Delta P=\left( \dfrac{P}{U\cos \varphi} \right)^2. R=\left( \dfrac{1000}{110. 0,9} \right)^2. 20=2040,6\ W$.
			\item Điện năng hao phí trên đường dây trong 30 ngày là: $A=\Delta P. t=2040,6. 30. 24=1469232\ Wh\approx 1469\ kWh$.
		\end{itemize}
	}
\end{vidu}
\loai{Tính độ sụt áp}
\begin{vidu} %[3P3-15.1K][loai1]
	Một xưởng sản xuất hoạt động đều đặn và liên tục 8 giờ mỗi ngày, 22 ngày trong một tháng. Điện năng lấy từ máy hạ áp có điện áp hiệu dụng ở cuộn thứ cấp là 220 V. Điện năng truyền đến xưởng trên một đường dây có điện trở tổng cộng là 0,08 $\Omega$. Trong một tháng, đồng hồ đo trong xưởng cho biết xưởng tiêu thụ 1900,8 số điện (1 số điện $= 1\ kWh$). Coi hệ số công suất của mạch luôn bằng 1. Độ sụt áp trên đường dây tải bằng
	\choice
	{\True 4 V}
	{1 V}
	{2 V}
	{8 V}
	\loigiai{
		\begin{itemize}\item Ta có: 
			\begin{align*}
				&P_\text{hạ thế}=\Delta P_{dây}+P_{nhà máy} \Leftrightarrow UI=I^2R+\dfrac{1900,8. 10^3. 3600}{8. 3600. 22}\\
				\Leftrightarrow\ &0,08I^2-220I+10800=0 \Rightarrow I=2700\ A\ \ct{(loại) hoặc}\ I=50\ A.
			\end{align*}
			\item Vậy độ sụt thế trên đường dây là: $\Delta U=IR=4\ V$.
		\end{itemize}
	}
\end{vidu}
\loai{Tính hiệu suất truyền tải}
\begin{vidu} %[3P3-15.1K][loai1] 
	Một trạm phát điện truyền đi với công suất 100 kW, điện trở đường dây tải điện là $8\ \Omega $. Điện áp ở hai đầu trạm là 1000 V. Nối hai cực của trạm với một biến thế có tỉ số vòng dây cuộn sơ cấp và cuộn thứ cấp $\dfrac{N_1}{N_2}=0,1$. Cho rằng hao phí trong máy biến áp không đáng kể, hệ số công suất máy biến áp bằng 1. Hiệu suất tải điện của trạm khi có máy biến áp là
	\choice
	{99\%}
	{90\%}
	{92\%}
	{\True 99,2\%}
	\loigiai{
	\begin{itemize}
			\item Nối cực của trạm phát điện với một biến thế có $k = 0.1\ \Rightarrow U_{\text{phát}} = 10000\ V$.
			\item Công suất hao phí được xác định bởi biểu thức: $\Delta P=R.\dfrac{P^2}{U^2}\cos \varphi =800\ W$. 
			\item Hiệu suất truyền tải điện năng là: $H=1-\dfrac{\Delta P}{P}=99,2\%$.
	\end{itemize}
}
\end{vidu}
\loai{Hiệu số chỉ công tơ điện}
\begin{vidu} %[3P3-15.1K][loai1] 
	Điện năng ở một trạm phát điện được truyền đi với công suất 200 kW. Hiệu số chỉ của các công tơ điện ở trạm phát và ở nơi thu sau mỗi ngày đêm chênh lệch nhau 480 kWh. Công suất điện hao phí trên đường dây tải điện là
	\choice
	{\True 20 kW}
	{40 kW}
	{83 kW}
	{100 kW}
	\loigiai{
		\begin{itemize}
			\item Điện năng hao phí sau mỗi ngày đêm là: $\Delta A=480\ kWh$.
			\item Công suất hao phí trên đường dây tải điện là: $\Delta P=\dfrac{\Delta A}{t}=\dfrac{480}{24}=20\ kW$.
		\end{itemize}
	}
\end{vidu}
\loai{Điện áp hiệu dụng đưa lên đường dây}
\begin{vidu} %[3P3-15.1K][loai1] 
	Từ một máy phát điện người ta muốn truyền tới nơi tiêu thụ bằng đường dây tải điện có điện trở 40 $\Omega$ và hệ số công suất bằng 1. Biết hiệu suất truyền tải là 98\% và nơi tiêu thụ nhận được công suất điện 196 kW. Điện áp hiệu dụng đưa lên đường dây là
	\choice
	{10 kV}
	{\True 20 kV}
	{40 kV}
	{30 kV}
	\loigiai{
		\begin{itemize}
			\item $\left\{ \begin{array}{l}
			H = \dfrac{{P'}}{P} \Rightarrow 0,98 = \dfrac{{196}}{P} \Rightarrow P = 200\ kW\\[0.3cm]
			\Delta P = \left( {1 - H} \right)P = 4\ kW
			\end{array} \right.$
			\item $\Delta P = \dfrac{{{P^2}}}{{{U^2}}}R \Rightarrow {4.10^3} = \dfrac{{{{\left( {{{200.10}^3}} \right)}^2}.40}}{{{U^2}}} \Rightarrow U = {20.10^3}\ V$. 
		\end{itemize}
	}
\end{vidu}
\loai{Tổng hợp}
\begin{vidu} %[3P3-15.4K][loai1] 
	Người ta truyền tải điện năng từ M đến N. Ở M dùng máy tăng thế và ở N dùng máy hạ thế, dây dẫn từ M đến N có điện trở $40\ \Omega $. Cường độ dòng điện trên dây là 50 A. Công suất hao phí trên dây bằng $5\%$ công suất tiêu thụ ở N và điện áp ở cuộn thứ cấp của máy hạ thế là 200 V. Biết dòng điện và điện áp luôn cùng pha và bỏ qua hao phí của máy biến thế. Tỉ số số vòng dây của máy hạ thế là
	\choice
	{100}
	{250}
	{\True 200}
	{20}
	\loigiai{
		\begin{itemize}
			\item Ta có: $I_1^2R=5\%.U_2I_2\cos \varphi_2\Rightarrow 50^2.40=5\%.200.I_2.1\Rightarrow I_2=10000\ A$.
			\item Mặt khác theo công thức máy biến áp:
			$\dfrac{N_2}{N_1}=\dfrac{U_2}{U_1}=\dfrac{I_1}{I_2\cos \varphi_2}=\dfrac{50}{10000.1}=0,005\Rightarrow \dfrac{N_1}{N_2}=200$.
		\end{itemize}
	}
\end{vidu}